\subsection{실험 조건 요약}

\label{sec:ch4_1}

KoBLEX의 한국어 법률 질문 226문항에서 ParSeR의 재정렬 모델을 Cross-Encoder(CE, \texttt{dragonkue/bge-reranker-v2-m3-ko})에서 Decision Model인 Jev(\texttt{jev-1.14.0})로 교체하여 평가하였다. 선택·답변 생성·자동 평가에는 \texttt{gpt-4o-2024-08-06}을 사용하였다. 두 조건은 예상 법령과 BM25 검색 후보를 공유하고, 예상 법령별 후보 100개를 재정렬한 후 상위 10개 중 LLM이 하나를 선택하였다. 후보 목록은 총 478개이며 선택 수를 정답 근거 수에 맞추지 않았다. CE 입력 형식(원 질문만 사용)은 평가 문항 30개에서 네 형식을 비교한 탐색적 진단으로 정하였으며, 공식 코드의 재정렬 입력과 같고 진단에서 CE에 가장 유리했던 이 형식을 Jev에도 적용하였다.

비교는 공개 코드 CE, 이를 수정한 CE, Jev의 세 조건으로 구성된다. 공개 코드는 현재 목록에서 선택한 번호를 첫 번째 목록에 적용하므로, 이를 보존한 조건은 재현 참조로 분리하였다. 앞의 두 조건 차이는 후보 목록 참조 수정 효과, 수정 CE와 Jev의 차이는 재정렬기 교체 효과이며, 이하 CE는 수정 조건을 뜻한다. 모든 비교는 평가(test) 문항의 결과를 이미 관찰한 뒤 수행한 탐색적 분석이다. 품질 차이의 95\% 신뢰구간(CI)은 출처나 정답 근거를 공유하는 151개 그룹의 대응 bootstrap 10,000회로 산출하였고, 다중 비교를 보정하지 않았으므로 가설 생성적 근거로만 해석한다(보정 결과는 S1). 문항 평균 점수(P·R·F1, Token F1, LF-Eval)는 0–100 척도의 점, 문항 비율(EM, 완전 회수율 등)은 \%와 \%p로 표기한다. 지연은 hop별 층화 무작위로 미리 추출한 60문항(1/2/3-hop 15/34/11)에서 별도로 측정하고, 151개 그룹 중 표본에 해당하는 52개로 같은 방식의 구간을 산출하였다. Jev는 원격 API로 다시 호출하면 상위 10개가 달라질 수 있으므로(4.3) Jev 품질 수치는 저장된 단일 실행의 값이며, CI는 문항 표본 변동만 반영하고 실행 간 변동은 포함하지 않는다.
